\section{Monopulse Angle Estimation}
Target angle sensing is one of the principal functions of radar systems. In early systems, scanning 
or sequential lobing was used to estimate angle, which involved mechanically steering the radar 
beam and comparing the signal strength across different beam positions. Monopulse systems on the
other hand use multiple beams simultaneously to estimate angle using the difference in amplitude
or phase across each of the received channels \cite{shermanbarton2011monopulse}.

In this thesis, we focus on amplitude comparison monopulse processing, which is commonly used in 
small radar systems due to its size and weight requirements. About a particular dimension, the angle 
of a target can be estimated with the equation, 
\begin{equation}
    \theta = f\left( \frac{S_1 - S_2}{S_1 + S_2} \right)
\end{equation}
where \(S_1\) and \(S_2\) are the signal strengths in the two channels being compared, and \(f\) 
is a function that maps the sum-difference ratio to an angle. This explanation is illustrated in 
Figure \ref{fig:background_monopulse_principle}. 

\begin{figure}
    \incfig[0.9]{background_monopulse_principle}
    \caption[Amplitude comparison monopulse angle estimation principle.]{Principle of amplitude comparison monopulse processing. The angle of the target is
    estimated as near the antenna axis, slightly toward receiver 1, because the amplitude from receiver 1
    is stronger than from receiver 2.}
    \label{fig:background_monopulse_principle}
\end{figure}

With a square pattern of four receiver channels (Figure \ref{fig:quadruplet}), azimuth estimations can be obtained by obtaining 
the sum-difference ratio of the left and right channels, and elevation estimations can be obtained
by obtaining the sum-difference ratio of the top and bottom channels. 
\begin{figure}
    \incfig[0.26]{tetrad}
    \caption[Four-channel monopulse system.]{A four-channel monopulse receiver pattern, also known as a quadruplet.
    Receivers are labelled \(A_{0,0}\), \(A_{0,1}\), \(A_{1,0}\) and \(A_{1,1}\) as shown.}
    \label{fig:quadruplet}
\end{figure}

\begin{equation}
    \delta_b = \frac{|a_{0,0} + a_{1,0}| - |a_{0,1} + a_{1,1}|}{|a_{0,0} + a_{1,0}| + |a_{0,1} + a_{1,1}|}
\end{equation}
\begin{equation}
    \delta_e = \frac{|a_{0,0} + a_{0,1}| - |a_{1,0} + a_{1,1}|}{|a_{0,0} + a_{0,1}| + |a_{1,0} + a_{1,1}|}
\end{equation}

Where \(a_{i,j}\) represents the signal strength from receiver \(A_{i,j}\). The relationship between 
the ratios and true angles is dependent on the specific radar system. 
Different receivers will have different gains. Calibration allows us to map the sum-difference ratios
to angle estimates in radians using two polynomial functions of the form,

\begin{equation}
    \theta_b = \sum_{i=0}^{4} \sum_{j=0}^{4-i} c_{bij} \delta_b^i \delta_e^j
    \qquad
    \theta_e = \sum_{i=0}^{4} \sum_{j=0}^{4-i} c_{eij} \delta_b^i \delta_e^j
\end{equation}

One for azimuth and one for elevation, where \(c_{bij}\) and \(c_{eij}\) are the coefficients 
obtained from calibration data.

With multiple quadruplets, multiple fields of view can be stitched together to produce a wide field of view, 
without the need to scan. Single receivers can be members of multiple quadruplets, resulting in an efficient 
use of hardware, as 
depicted in Figure \ref{fig:overlapping_quadruplets}. 
\begin{figure}
    \incfig[0.45]{overlapping_tetrads}
    \caption[Overlapping quadruplets.]{Overlapping quadruplets allow for a wide field of view, a higher capacity
    to track multiple targets, and efficient use of hardware, as a single receiver can be a member of 
    multiple quadruplets.}
    \label{fig:overlapping_quadruplets}
\end{figure}